   \documentclass{article}
    
    \textwidth=6.5in
    \textheight=8.6in
    \voffset=-54pt
    \oddsidemargin=0in
    \evensidemargin=0in
    \setlength{\parskip}{8pt}
    \setlength{\parindent}{0pt}
    
  
    \usepackage[normalem]{ulem}
    \usepackage{microtype}
    
    
    \usepackage{scalerel,amssymb}
\def\mcirc{\mathbin{\scalerel*{\circ}{j}}}
\def\msquare{\mathord{\scalerel*{\Box}{gX}}}

    \usepackage{diffcoeff}

    \usepackage{amsmath, amssymb}
    \usepackage{ifthen}
    
    \usepackage[inline]{enumitem}
    \setlist{topsep=0pt}
    
    \usepackage{xcolor}

    \newcommand{\eqdots}{\setbox0=\hbox{=}%
       \mathrel{\hbox to\wd0{\hss\vdots\hss}}}
    \usepackage[percent]{overpic}
    \usepackage{pifont}
   \usepackage{siunitx} 
    
    % change hyperref options
    \PassOptionsToPackage{hyphens}{url}

\usepackage[colorlinks,linkcolor=blue,urlcolor=blue,breaklinks]{hyperref}

    
    \usepackage{graphicx}
    \usepackage{multicol}
    \usepackage{framed}
    \usepackage{array}
    \usepackage{soul}
    \usepackage{wrapfig}
    
    \usepackage{pgfplots}
    \pgfplotscreateplotcyclelist{mycolor}{%
      black, blue, red, green!70!black,magenta,
      purple, cyan, orange
    }
    \pgfplotsset{
      cycle list name=mycolor,
      axis lines=middle,
      every axis plot/.style={no marks, thick},
      every axis/.style={
        enlarge x limits=false,
        enlarge y limits=auto,
        xlabel style={at={(current axis.right of origin)},anchor=west},
        ylabel style={at={(current axis.above origin)},anchor=south},
        % extra x and y ticks for asymptotes
        extra x tick style={grid=major, grid style={dashed,black}},
        extra y tick style={grid=major, grid style={dashed,black}},
        /pgf/number format/1000 sep={},
      },
      tick label style = {font=\footnotesize, fill=white},
      scaled ticks=false,
      scaled y ticks=false,
      unbounded coords=jump,
      samples=101,
      scatter/classes={
        open={mark=*,draw=black,fill=white, mark size=2, thin},
        closed={mark=*,draw=black,fill=black, mark size=2, thin}
      },
    }
      
    
    \usepackage{tikz}
    \usetikzlibrary{
      % enables the snake paths
      decorations.pathmorphing,%
      % enables bigger arrows
      decorations.markings,%
      decorations.pathreplacing,%
      decorations.text,%
      calc,%
      patterns,%
      shapes.geometric,%
      arrows,
      decorations.shapes,
      positioning,
      plotmarks, angles, quotes,
      pgfplots.groupplots
    }
    \tikzset{>=latex} % sets default arrow type in tikz
    \tikzset{font=\small}
    
      
    \usepackage{multirow, bigdelim}
    
    % change to \usepackage[sols]{handout} to see solutions
    \usepackage[]{handout}
        
    \newcommand{\lheq}{\stackrel{\text{L'H}}{=}}
\newcommand{\nolheq}{\mathrel{\makebox[\widthof{$\lheq$}]{=}}}

    \definecolor{skyblue}{rgb}{0, 0.5, 1}
    \usepackage{tcolorbox}
\newenvironment{feedback}{%
  \begin{tcolorbox}[colframe=skyblue, colback=skyblue!10!white]
  }{\end{tcolorbox}}
      
    
    \newcommand\iscurrentenumi[1]{%
      \ifthenelse{\equal{#1}{\theenumi}}{}{#1}}
    \setenumerate[1]{ref=\#\arabic*}
    \setenumerate[2]{ref=\protect\iscurrentenumi{\theenumi}(\alph*)}

    
      
    
    \newlength{\tipmargin}
    \makeatletter
    
    \makeatother
    
      
    \ifthenelse{\isundefined{\C}}{%
      \newcommand{\C}{\mathbb{C}}%
    }{\renewcommand{\C}{\mathbb{C}}}
    \ifthenelse{\isundefined{\R}}{%
      \newcommand{\R}{\mathbb{R}}%
    }{\renewcommand{\R}{\mathbb{R}}}
    
    \definecolor{purple}{rgb}{0.5,0,1.0}
    \pgfplotsset{holdot/.style={color=red,fill=white,only marks,mark=*}}
    \pgfplotsset{soldot/.style={color=black,fill=black,only marks,mark=*}}
\pgfplotsset{compat=1.13}
\usetikzlibrary{quotes,arrows.meta}

    \pgfplotsset{holdot/.style={color=red,fill=white,only marks,mark=*}}

\newcommand{\answerbox}[3][]{
    \fbox{\begin{minipage}[t][#2][c]{#3}
        Final answer: #1
    \end{minipage}}
    }

\newcommand{\answerboxd}[3][]{
    \fbox{\begin{minipage}[t][#2][c]{#3}
        #1
    \end{minipage}}
    }
    
\begin{document}

 

 \noindent
{\large \mbox{} \hfill \bf Name:\underline{\hspace{2.8in}}}

 {\mbox{} \hfill \it (Please write your name neatly in print.)}


 
\medskip
 
\noindent
{\large\bf
Practice Trig Skills Check 2 \\ Math Mb Spring 2025}
\noindent






\iffalse
\begin{center}
\begin{tabular}{|c|c|c|}
\hline
Problem & Points & Score\\ \hline
& &\\
1 &  & \\ 
\hline
&& \\
2&  &  \\ 
\hline
& &\\
3&&  \\ 
 \hline
& &\\
Total&  &  \\ 
& &\\ \hline

\end{tabular}
\end{center}
\fi


%{\large \em Please circle your section:}


%\begin{center}
%{\small
%\begin{tabular}{ccccccccc}
%Erica Dinkins & Matt Cavallo &  Shanelle Davis  & Robin Gottlieb & 	Mike Koske & Hakim Walker \\
%MWF 9am  & MWF 9am &  MWF 10:30am  &MWF 10:30am & MWF 12pm   & MWF 12pm\\
%\\\\
% &&  &     \\
% &&  &
%\end{tabular}

%}
%\vspace{ -.5in} 
%{\small
%\begin{tabular}{cccccccc}
 %Hannah Constantin & Matt Cavallo & 	Justin Hancock  & Grant Barkley \\
 %MWF 12pm &  MWF 12pm  & MWF 1:30pm  & MWF 3pm  \\
%\\\\
 %&&       \\
 %&&  
%\end{tabular}


%}
%\end{center}


{\Large
\begin{center}\textbf{Directions---Please read carefully!} \end{center}}

This is a practice test. The same directions will appear on the actual Skills Check, so read them carefully.

\hrulefill 

You have 60 minutes to complete this Skills Check. \textbf{Calculators or any other aids are not permitted.}  


For each problem, write your final answer in the answer box provided. You may do work on any page and may ask for scratch paper, but \textbf{no work or writing outside of the answer box will be graded.}

Write neatly, as answers that cannot be read will not receive credit. 

   

\begin{center}\textbf{Good luck!}\end{center}


\newpage

\begin{enumerate}

\item Evaluate each expression.

    \begin{multicols}{2}
    \begin{enumerate}[label=(\alph*)]
        \item $\cos(7\pi/4)$
        \item $\sin(\pi/2)$
    \end{enumerate}
    \end{multicols}

    \vfill 

        \begin{multicols}{2}
                \answerbox[\solonly{$\frac{\sqrt{2}}{2}$}]{0.8in}{2.8in}
                \answerbox[\solonly{$1$}]{0.8in}{2.8in}
        \end{multicols}

        \begin{multicols}{2}
    \begin{enumerate}[start=3,label=(\alph*)]
        \item $\cos(5\pi/6)$
        \item $\sin(19\pi/6)$
    \end{enumerate}
    \end{multicols}

    \vfill

        \begin{multicols}{2}
                \answerbox[\solonly{$-\frac{\sqrt{3}}{2}$}]{0.8in}{2.8in}
                \answerbox[\solonly{$-\frac{1}{2}$}]{0.8in}{2.8in}
        \end{multicols}

    \begin{multicols}{2}
    \begin{enumerate}[start=5,label=(\alph*)]
        \item $\cos(-4\pi/3)$
        \item $\sin(17\pi/4)$
    \end{enumerate}
    \end{multicols}

    \vfill 

        \begin{multicols}{2}
                \answerbox[\solonly{$-\frac{1}{2}$}]{0.8in}{2.8in}
                \answerbox[\solonly{$\frac{\sqrt{2}}{2}$}]{0.8in}{2.8in}
        \end{multicols}



\newpage 

\item Solve the equation $\cos(\theta)=\frac{1}{2}$ on the interval $[0,2\pi)$.

\vfill

\answerbox[\solonly{$\theta=\dfrac{\pi}{3}, \dfrac{5\pi}{3}$}]{0.8in}{5.5in}

 \item $A$ is an angle such that $\cos A=-\frac{1}{3}$ and $\sin A < 0$. Find $\sin A$.

        \vfill

        \answerbox[\solonly{$\sin A = -\dfrac{2\sqrt{2}}{3}$}]{0.8in}{5.5in}

\item A particle is in a circular particle accelerator with a radius of 100 m. In one millisecond, the particle's path makes an arc that subtends an angle of $\pi/4$ radians at the center of the circle. What is the distance traveled by the particle?

\vfill

\answerbox[\solonly{$25\pi$ meters}]{0.8in}{5.5in}


\newpage 

\item
Consider $f(t) =1 -3 \sin(\pi t)$. 
    \begin{enumerate}
        \item Find the midline, amplitude, and period.

        \vfill

    \answerboxd[Midline: $y=\solonly{1}$]{.5 in}{5.5 in}

\vspace{.1in}
    
\answerboxd[Amplitude $=\solonly{3}$]{.5 in}{5.5 in}

\vspace{.1in}

\answerboxd[Period $=\solonly{2}$]{.5 in}{5.5 in}

\vspace{.1in}

    \item Which of these is the graph of $y=f(t)$? Circle the letter above your choice.
    \begin{center}
\begin{tikzpicture}
    \begin{groupplot}[
        group style = {group size=2 by 2},
        grid = none,
        xtick={0, 2, 4, 6},
        xticklabels={0,2,4, 6},
        ytick={-3,-2,-1, 0, 1, 2, 3, 4},
        yticklabels={-3, -2, -1, 0, 1, 2, 3, 4},
        width = 2.3in,
        xmin=-1, xmax=6, ymin=-3, ymax=4,
        restrict y to domain=-100:100]
        
    \nextgroupplot[title=(A)]
    \addplot[draw=none] coordinates {(0,0)};
    \addplot[domain = -1:8, color = black]{-3*sin(deg(x)};
   

    \nextgroupplot[title=(B)]
    \addplot[draw=none] coordinates {(0,0)};
    \addplot[domain = -1:8, color = black]{3* sin(deg(pi*x))+1};


    \probonly{\nextgroupplot[title=(C)]}
    \solonly{\nextgroupplot[title=\fbox{(C)}]}
    \addplot[draw=none] coordinates {(0,0)};
    \addplot[domain = -1:8, color = black]{-3*sin(deg(pi*x))+1};

    \nextgroupplot[title=(D)]
    \addplot[draw=none] coordinates {(0,0)};
    \addplot[domain = -1:8, color = black]{3*cos(pi*deg(x)};

    \end{groupplot}
\end{tikzpicture}
\end{center}
    \end{enumerate}

\newpage


    % \item Determine the midline, amplitude, and period of $y=5-\cos(\pi t)$.

    %     \vfill

    %     \answerboxd[Midline: $y=\cfitb{0.75in}{5}$\hspace{12pt} Amplitude $=\cfitb{0.75in}{1}$ \hspace{12pt} Period $=\cfitb{0.75in}{2}$]{0.8in}{5.5in}

    
    % \item Write an equation for the sinusoidal graph shown.

    % \begin{tikzpicture}
    %     \begin{axis}[
    %     xlabel=$t$,
    %     ylabel=$y$,
    %     samples=100,
    %     ytick=\empty, xtick=\empty,
    %     ylabel style={rotate=-90}
    %     ]
    %         \addplot[domain=-30:30] {3*sin(deg(x/4))-1};
    %         \addplot[soldot] coordinates {(0,-1) (6.28,2)};
    %         \node[right] at (0,-1) {$(0,-1)$};
    %         \node[above] at (6.28,2) {$(2\pi,2)$};
    %     \end{axis}
    % \end{tikzpicture}

    % \vspace{\stretch{0.5}}

    %     \answerbox[\solonly{$y=3\sin(\frac{1}{4}t)-1$}]{0.8in}{5.5in}
    
\newpage

\item Determine whether each statement is an identity (true for all values). Write ``True" or ``False."


    \begin{multicols}{2}
            \begin{enumerate}
        \item  $\sin(-\theta)=-\sin(\theta)$

        \item $\cos(A+B)=\cos(A)+\cos(B)$

    \end{enumerate}
    \end{multicols}

    \vfill 

        \begin{multicols}{2}
                \answerbox[\solonly{True}]{0.8in}{2.5in}
                
                \answerbox[\solonly{False}]{0.8in}{2.5in}
        
    \end{multicols}

    \begin{multicols}{2}
            \begin{enumerate}[start=3]
        \item $\sin^2\theta = \cos^2\theta-1$
        \item $\sin(\theta-\frac{\pi}{2})=\cos(\theta)$
    \end{enumerate}
    \end{multicols}
\vfill 

    \begin{multicols}{2}
                \answerbox[\solonly{False}]{0.8in}{2.5in}
                
                \answerbox[\solonly{False}]{0.8in}{2.5in}
        
    \end{multicols}


\newpage 
        
\item Evaluate each expression. If the expression is undefined, write ``undefined."

    \begin{multicols}{2}
    \begin{enumerate}[label=(\alph*)]
        \item $\tan(3\pi/4)$
        \item $\tan(7\pi/6)$
    \end{enumerate}
    \end{multicols}

    \vfill 

        \begin{multicols}{2}
                \answerbox[\solonly{$-1$}]{0.8in}{2.8in}
                \answerbox[\solonly{$\dfrac{\sqrt{3}}{3}$}]{0.8in}{2.8in}
        \end{multicols}

        \begin{multicols}{2}
    \begin{enumerate}[start=3,label=(\alph*)]
        \item $\sec(2\pi/3)$
        \item $\csc(5\pi/6)$
    \end{enumerate}
    \end{multicols}

    \vfill

        \begin{multicols}{2}
                \answerbox[\solonly{$-2$}]{0.8in}{2.8in}
                \answerbox[\solonly{$2$}]{0.8in}{2.8in}
        \end{multicols}

    \begin{multicols}{2}
    \begin{enumerate}[start=5,label=(\alph*)]
        \item $\cot(\pi)$
        \item $\tan(\pi)$
    \end{enumerate}
    \end{multicols}

    \vfill 

        \begin{multicols}{2}
                \answerbox[\solonly{undefined}]{0.8in}{2.8in}
                \answerbox[\solonly{$0$}]{0.8in}{2.8in}
        \end{multicols}

\newpage

\item At a distance of 1000 feet, a surveyor on the ground measures the angle of elevation to the top of the Empire State Building to be \ang{55.5}. Using this information, find the height of the Empire State Building. You may leave unevaluated trigonometric expressions in your answer. Include units.

\vfill

\answerbox[\solonly{$1000\tan(\ang{55.5})\text{ ft}$}]{0.8in}{5.5in}


\item Sarah is flying a kite and has tied it to the ground. The 200-ft string attached to the kite makes an angle of \ang{43} with the ground. How high above the ground is the kite flying? You may leave unevaluated trigonometric expressions in your answer. Include units.

\vfill

\answerbox[\solonly{$200\sin(\ang{43})\text{ ft}$}]{0.8in}{5.5in}

% \item A passenger in an airplane at an altitude of 10 kilometers looks out the window and sees a town to the east of the plane. The angle of depression to the town is \ang{22}. How many kilometers east of the plane is the town? You may leave unevaluated trigonometric expressions in your answer.

        
\end{enumerate}



\end{document}